% =========================================================
%   CAPÍTULO 1: INTRODUCCIÓN
% =========================================================
\chapter{Introducción}
\label{cap:introduccion}

% =========================================================
\section{Planteamiento del problema}
\label{sec:planteamiento}
% =========================================================

El Laboratorio de Fotofísica y Películas Delgadas del ICAT cuenta con
un arreglo fotoacústico pulsado destinado a la caracterización de
líquidos de baja absorción óptica. El arreglo es funcional y ha
sostenido trabajos experimentales previos, pero su operación es
enteramente manual: cada punto experimental exige coordinar el disparo
del láser, la adquisición del osciloscopio y el registro de la
temperatura de la muestra, y anotar aparte las condiciones bajo las
cuales se tomó la medición. Los fundamentos físicos que justifican esa
configuración se desarrollan en el \cref{cap:fundamentos}.

La consecuencia más visible de esa operación manual es la lentitud, y
es también la menos importante. El problema de fondo es la
correspondencia entre un dato y las condiciones que lo produjeron. Una
captura guardada a mano queda descrita por lo que el operador alcanzó a
anotar, y esa anotación vive en un cuaderno, separada del archivo: puede
extraviarse, puede contradecir al archivo, y no existe forma de detectar
la discrepancia meses después, cuando el dato se analiza. Un experimento
cuyo interés reside en comparar señales tomadas bajo condiciones
distintas —niveles de energía, temperaturas, líquidos— descansa por
completo sobre esa correspondencia.

A ello se añade que el láser es un recurso compartido del laboratorio,
de modo que el arreglo se desmonta y se restablece entre sesiones. La
geometría efectiva cambia, y con ella parámetros que condicionan la
interpretación de la señal. Sin un registro que acompañe a cada captura,
la reproducibilidad entre sesiones no puede afirmarse ni desmentirse.

% =========================================================
\section{Premisa de trabajo}
\label{sec:premisa}
% =========================================================

El trabajo parte de la siguiente premisa:

\begin{quote}
Un sistema que registre cada captura junto con las condiciones
instrumentales que la produjeron permite reconstruir a posteriori el
estado del arreglo en el momento de la adquisición, incluidas
condiciones anómalas que no fueron observables durante la sesión.
\end{quote}

Se enuncia como premisa y no como hipótesis porque no se contrasta
contra la naturaleza sino contra el desempeño de un artefacto
construido: su comprobación consiste en mostrar que el registro
conservado basta para reconstruir situaciones que en su momento pasaron
inadvertidas. El \cref{cap:resultados} documenta cinco anomalías
ocurridas durante el desarrollo, cuatro de las cuales no eran
observables mientras transcurrían, y todas ellas caracterizables meses
después a partir de los archivos de sesión.

% =========================================================
\section{Objetivos}
\label{sec:objetivos}
% =========================================================

\subsection*{Objetivo general}

Desarrollar un sistema de automatización que integre el láser de
excitación, el osciloscopio de adquisición y la medición de temperatura
de la muestra bajo una sola aplicación, ejecute secuencias de medición
sin intervención del operador entre puntos, y conserve para cada
captura los metadatos necesarios para reconstruirla en unidades físicas
y para conocer las condiciones bajo las cuales fue adquirida.

\subsection*{Objetivos particulares}

\begin{enumerate}[label=\arabic*.]
  \item Implementar el control programático del láser a través de la
        biblioteca de enlace dinámico del fabricante, incluido un
        procedimiento de retorno a estado seguro ante la conclusión de
        una secuencia, una parada de emergencia o un fallo de
        comunicación.

  \item Implementar el control del osciloscopio sobre el protocolo
        VXI-11, con transferencia binaria de la forma de onda y lectura
        del preámbulo que permite su conversión posterior a unidades
        físicas.

  \item Definir e implementar la interfaz de adquisición de temperatura
        del sistema: el contrato de comunicación que la aplicación
        espera, la distinción entre una lectura válida y la pérdida de
        un sensor, y la correspondencia estable entre cada valor
        recibido y el dispositivo que lo produjo.

  \item Implementar los dos modos de secuencia automática previstos para
        el experimento: adquisición durante un intervalo de tiempo fijo
        y adquisición gobernada por la temperatura de la muestra durante
        su enfriamiento.

  \item Definir un formato de almacenamiento de sesión que conserve las
        muestras crudas del instrumento, los parámetros vigentes de cada
        dispositivo al momento de la captura y una marca explícita de
        las condiciones anómalas detectadas durante la adquisición.

  \item Validar la lógica del sistema sin acceso a los instrumentos,
        mediante simulación de perfiles térmicos, recorrido completo de
        la cadena de medición y almacenamiento, y sesiones sintéticas
        construidas para provocar cada una de las condiciones que la
        herramienta de verificación debe detectar.
\end{enumerate}

% =========================================================
\section{Alcance y delimitación}
\label{sec:alcance}
% =========================================================

Conviene establecer desde ahora qué queda fuera de este trabajo, pues
varias de las limitaciones que se discuten en los capítulos siguientes
son consecuencia de esta delimitación y no defectos del sistema
desarrollado.

El trabajo desarrolla el sistema de automatización de un arreglo
experimental preexistente. No modifica la celda fotoacústica, diseñada y
manufacturada en un trabajo anterior, ni el hidrófono, cuyo ancho de
banda condiciona la forma de las señales registradas; ambos se describen
en el \cref{cap:arreglo} como condiciones de frontera del experimento.
Tampoco incorpora un sensor calibrado de energía del pulso láser, de
modo que las amplitudes adquiridas no son normalizables y el sistema, por
completo que sea su registro, no sostiene por sí solo afirmaciones
cuantitativas sobre el coeficiente de absorción del medio.

El trabajo tampoco caracteriza líquidos. Las sesiones experimentales
reportadas se realizaron con agua y su propósito es evaluar la
herramienta —su repetibilidad, su capacidad de distinguir un cambio
deliberado de parámetros y la integridad de su registro—, no las
propiedades del líquido empleado. La caracterización de líquidos es lo
que el sistema hace posible, no lo que este trabajo entrega.

% =========================================================
\section{Metodología}
\label{sec:metodologia}
% =========================================================

El desarrollo se organizó en cuatro etapas.

La primera consistió en resolver el control de cada instrumento por
separado, atendiendo a las particularidades de su interfaz: la carga de
una biblioteca de 64 bits que resuelve sus dependencias por rutas
relativas en el caso del láser, y el orden de configuración de la
transferencia binaria en el caso del osciloscopio. La segunda integró
los tres subsistemas bajo una sola aplicación de escritorio, con las
operaciones de hardware ejecutándose en hilos independientes del de la
interfaz gráfica y comunicando sus resultados mediante señales.

La tercera etapa fue la validación sin instrumentos. El acceso al láser
y al osciloscopio se limita a las sesiones reservadas, y una parte de
las condiciones que interesa verificar —la pérdida de un dispositivo a
mitad de secuencia, el rebote de una lectura alrededor del umbral de
disparo, la escritura de un registro incompleto— no puede provocarse a
voluntad frente al equipo sin arriesgar la muestra o el instrumento. Se
adoptó por ello el criterio de validar por simulación y sesiones
sintéticas todo aquello que depende de la lógica del programa,
reservando para el laboratorio únicamente lo que depende de la respuesta
efectiva del equipo. La herramienta de verificación de datos se sometió
además a una prueba de mutación, en la que cada una de sus
comprobaciones se desactivó por separado para confirmar que su ausencia
hacía fallar el caso correspondiente.

La cuarta etapa evaluó el comportamiento del sistema integrado en
sesiones reales de laboratorio, y documentó las anomalías observadas
durante el desarrollo junto con el nivel al que fue posible detectarlas.

% =========================================================
\section{Estructura del documento}
\label{sec:estructura}
% =========================================================

El \cref{cap:fundamentos} presenta los fundamentos del efecto
fotoacústico, su manifestación en líquidos de baja absorción, la
necesidad de caracterizar este tipo de medios y las técnicas
disponibles para hacerlo. El \cref{cap:arreglo} describe el arreglo
experimental: la celda, los equipos de medición y la trayectoria óptica
heredados del laboratorio, y los tres subsistemas de instrumentación
automatizados en este trabajo. El \cref{cap:resultados} presenta el
sistema desarrollado, la evidencia que sostiene lo que cada subsistema
declara, el comportamiento del conjunto en sesiones reales y los modos
de falla observados. El \cref{cap:conclusiones} recoge el balance del
trabajo y las líneas de continuación identificadas durante su
desarrollo. El \cref{ap:tiempos_arribo} desarrolla el cálculo de los
tiempos de arribo acústicos en la celda y el procedimiento para
recalcularlos cuando la geometría del montaje difiere de la nominal.
